\documentclass[11pt,a4paper]{article}

% ---------- Packages ----------
\usepackage[utf8]{inputenc}
\usepackage[T1]{fontenc}
\usepackage{amsmath,amssymb,amsthm}
\usepackage{mathtools}
\usepackage{empheq}
\usepackage{bm}
\usepackage{geometry}
\usepackage{booktabs}
\usepackage{array}
\usepackage{siunitx}
\usepackage{hyperref}
\usepackage{xcolor}
\usepackage{enumitem}

\geometry{margin=1in}
\hypersetup{
    colorlinks=true,
    linkcolor=blue!50!black,
    urlcolor=blue!50!black,
    citecolor=blue!50!black
}

% ---------- Macros ----------
\newcommand{\OmegaT}{\Omega_t}
\newcommand{\Rt}{R_t}
\newcommand{\At}{A_t}
\newcommand{\Tt}{T_t}
\newcommand{\thetat}{\theta_t}
\newcommand{\vi}{v_i}
\newcommand{\viss}{v_{i,\text{ss}}}
\newcommand{\bvi}{\bar{v}_i}
\newcommand{\bT}{\bar{T}}
\newcommand{\btheta}{\bar{\theta}}
\newcommand{\blambda}{\bar{\lambda}}
\newcommand{\bCT}{\bar{C}_T}
\newcommand{\Tcmd}{T_t^{\text{cmd}}}
\newcommand{\CTcmd}{C_T^{\text{cmd}}}
\newcommand{\lcmd}{\lambda^{\text{cmd}}}

\title{\textbf{Tail Rotor Dynamics for an RC Helicopter}\\[4pt]
\large Thrust, Induced Inflow, and Nonlinear Lead Compensation}
\author{}
\date{}

\begin{document}
\maketitle

\begin{abstract}
\noindent This note develops a control-oriented model of an RC helicopter tail rotor.
Starting from blade-element momentum theory (BEMT), we derive the thrust-versus-pitch
relation, add dynamic inflow via the Pitt--Peters first-order model, invert the
coupled system to obtain the blade pitch required to track a commanded thrust, and
then extend the inverse to handle the static and dynamic nonlinearities of inflow
softening. The result is a nonlinear lead compensator suitable for a tail-rotor
inner loop. Realistic parameter ranges for RC tail rotors are tabulated.
\end{abstract}

\tableofcontents
\vspace{1em}
\hrule
\vspace{1em}

% =====================================================================
\section{Tail Rotor Thrust vs.\ Blade Pitch}
% =====================================================================

For a tail rotor with collective pitch \(\thetat\) (no cyclic), we use
\emph{Blade Element Momentum Theory} (BEMT).

\subsection{Blade Element Theory}

Consider a blade element at radial station \(r\) on a blade of chord \(c\),
rotating at \(\OmegaT\). The local flow has
\begin{align}
U_T &= \OmegaT\,r \qquad &&\text{(tangential)},\\
U_P &= \vi + \mu_z\,\OmegaT\,\Rt \qquad &&\text{(perpendicular)},
\end{align}
where \(\vi\) is the induced velocity and \(\mu_z\) is the axial inflow from
fuselage translation (sideways for a tail rotor).

The inflow angle is
\begin{equation}
\phi = \tan^{-1}\!\left(\frac{U_P}{U_T}\right) \approx \frac{U_P}{U_T},
\end{equation}
and the effective angle of attack of the blade element is
\(\alpha = \thetat - \phi\).
Assuming a linear lift curve, \(C_\ell = a\,\alpha\) with
\(a\approx \SI{5.7}{\per\radian}\), the elemental thrust per blade is
\begin{equation}
dT = \tfrac{1}{2}\rho\,(U_T^2+U_P^2)\,c\,C_\ell\,dr
   \approx \tfrac{1}{2}\rho\,(\OmegaT r)^2\,c\,a\,(\thetat-\phi)\,dr.
\end{equation}

\subsection{Integrated Thrust (Uniform Inflow)}

Using the non-dimensional inflow ratio
\(\lambda = (\vi + V_\infty)/(\OmegaT\Rt)\) and integrating over the blade with
\(N_b\) blades:
\begin{empheq}[box=\fbox]{equation}
C_T \;=\; \frac{\sigma a}{2}\!\left(\frac{\thetat}{3} - \frac{\lambda}{2}\right),
\label{eq:CT_BEMT}
\end{empheq}
where \(\sigma = N_b c/(\pi\Rt)\) is the rotor solidity. The dimensional thrust is
\begin{equation}
\Tt = C_T\,\rho\,\At\,(\OmegaT\Rt)^2, \qquad \At = \pi\Rt^2.
\end{equation}

For a small RC tail rotor at constant \(\OmegaT\), this is almost linear in \(\thetat\):
\begin{equation}
\Tt \approx K_1\,\thetat - K_2\,\vi,
\label{eq:linT}
\end{equation}
with
\begin{equation}
K_1 = \frac{\rho\At(\OmegaT\Rt)^2\,\sigma a}{6},
\qquad
K_2 = \frac{\rho\At\,\OmegaT\Rt\,\sigma a}{4}.
\label{eq:K1K2}
\end{equation}

% =====================================================================
\section{Induced Inflow: Momentum Theory}
% =====================================================================

\subsection{Steady Momentum Balance}

Momentum theory at the rotor disk equates thrust to the rate of change of
fluid momentum:
\begin{equation}
\Tt = 2\rho\At\,\vi\sqrt{(V_\infty+\vi)^2 + V_\perp^2}.
\label{eq:momentum}
\end{equation}
In hover (\(V_\infty=V_\perp=0\)) this reduces to
\begin{equation}
v_h = \sqrt{\frac{\Tt}{2\rho\At}}.
\label{eq:vh}
\end{equation}
In non-dimensional form, \(\lambda_i = C_T/\bigl(2\sqrt{\mu^2+\lambda^2}\bigr)\).
Combining with~\eqref{eq:CT_BEMT} yields the implicit equation
\begin{equation}
\lambda_i \;=\; \frac{\sigma a}{16}\!\left[\sqrt{1+\frac{64\thetat}{3\sigma a}\sqrt{\mu^2+\lambda^2}}-1\right]
\quad\text{(hover, }\mu=0\text{)}.
\label{eq:lambda_implicit}
\end{equation}

\subsection{Dynamic Inflow (Pitt--Peters)}

Steady momentum theory ignores the finite time taken by the wake to develop.
For control design we model the uniform-inflow component as a state:
\begin{empheq}[box=\fbox]{equation}
\tau_i\,\dot{\vi} + \vi \;=\; \viss,
\label{eq:pittpeters}
\end{empheq}
where \(\viss\) is the steady-state induced velocity from~\eqref{eq:momentum}.
A common closed form for the time constant is
\begin{equation}
\tau_i \;=\; \frac{0.849}{4\,\OmegaT}\cdot\frac{1}{\lambda_i^{\text{trim}}}
       \;=\; \frac{0.849\,\Rt}{4\,v_h}.
\label{eq:tau_i}
\end{equation}
The factor \(0.849\) comes from the apparent-mass coefficient \(8/(3\pi)\) of
an impermeable disk in the Pitt--Peters derivation.

\subsection{Coupled State-Space Form}

Combining BEMT with dynamic inflow:
\begin{empheq}[box=\fbox]{align}
\Tt &= \rho\At(\OmegaT\Rt)^2\,\frac{\sigma a}{2}\!\left(\frac{\thetat}{3}-\frac{\vi+V_\infty}{2\OmegaT\Rt}\right),\\[4pt]
\dot{\vi} &= \frac{1}{\tau_i}\!\left(\frac{\Tt}{2\rho\At\sqrt{(V_\infty+\vi)^2+V_\perp^2}} - \vi\right).
\end{empheq}

\paragraph{Full Pitt--Peters (three-state).}
For higher fidelity, Pitt--Peters extends this to three states---uniform inflow
\(\lambda_0\) plus longitudinal and lateral harmonics \(\lambda_s,\lambda_c\):
\begin{equation}
[M]\begin{Bmatrix}\dot{\lambda}_0\\\dot{\lambda}_s\\\dot{\lambda}_c\end{Bmatrix}
+ [L]^{-1}\begin{Bmatrix}\lambda_0\\\lambda_s\\\lambda_c\end{Bmatrix}
= \begin{Bmatrix}C_T\\ C_{Mx}\\ C_{My}\end{Bmatrix},
\end{equation}
with \([M]=\mathrm{diag}\!\bigl(\tfrac{8}{3\pi},\tfrac{16}{45\pi},\tfrac{16}{45\pi}\bigr)\)
and \([L]\) depending on advance ratio and wake skew. For a small tail rotor in
near-hover the one-state version is almost always sufficient.

% =====================================================================
\section{Inverse Problem: Blade Pitch from Commanded Thrust}
% =====================================================================

Given a desired thrust \(\Tcmd(t)\), find the pitch \(\thetat(t)\).
The blade-element relation has no dynamics; all time-dependence enters
through \(\vi\).

\subsection{Instantaneous Inverse}

Solving~\eqref{eq:linT} for \(\thetat\):
\begin{empheq}[box=\fbox]{equation}
\thetat(t) \;=\; \frac{\Tcmd(t) + K_2\bigl(\vi(t)+V_\infty(t)\bigr)}{K_1}.
\label{eq:inv_instant}
\end{empheq}
This is exact at every instant \emph{given the current} \(\vi\).
The inflow state must still be propagated via~\eqref{eq:pittpeters}.

\subsection{Linearised Transfer Function (Hover)}

Linearise about a hover trim \((\bT,\bvi,\btheta)\). The hover relation
\(\bT=2\rho\At\bvi^{\,2}\) gives
\begin{equation}
\delta\viss = \frac{1}{4\rho\At\bvi}\,\delta T \;\equiv\; k_v\,\delta T.
\label{eq:kv}
\end{equation}
The Pitt--Peters dynamics become \((\tau_i s + 1)\delta\vi = k_v\delta T\),
and the linearised BEMT relation is \(\delta T = K_1\delta\thetat - K_2\delta\vi\).
Eliminating \(\delta\vi\):
\begin{empheq}[box=\fbox]{equation}
\frac{\delta T(s)}{\delta\thetat(s)} \;=\; \frac{K_1(\tau_i s + 1)}{\tau_i s + (1+K_2k_v)}.
\label{eq:plant_tf}
\end{empheq}
This is a lead--lag response: a thrust spike of \(K_1\delta\thetat\) immediately,
settling at the lower DC gain \(K_1/(1+K_2k_v)\) once inflow builds up.

The inverse---a lead compensator---is
\begin{empheq}[box=\fbox]{equation}
\frac{\delta\thetat(s)}{\delta T(s)} \;=\; \frac{\tau_i s + (1+K_2 k_v)}{K_1(\tau_i s + 1)}.
\label{eq:comp_tf}
\end{empheq}

\paragraph{Interpretation of \(k_v\).} \(k_v\) is the linearised sensitivity of
steady-state induced velocity to thrust at trim:
\begin{equation}
k_v = \frac{\partial \viss}{\partial T}\bigg|_{\text{trim}} = \frac{1}{4\rho\At\bvi} = \frac{\bvi}{2\bT}.
\end{equation}
Units: \si{\metre\per\second\per\newton}.
The product \(K_2 k_v\) is dimensionless and captures the fractional thrust loss
due to inflow build-up at low frequency:
\begin{equation}
\frac{\text{steady gain}}{\text{instantaneous gain}} = \frac{1}{1+K_2 k_v}.
\end{equation}

\paragraph{Self-consistent form.} Computing \(K_2\) and \(k_v\) separately and
multiplying \emph{overestimates} the coupling, because both depend on \(\bvi\)
through the trim condition. The self-consistent dimensionless coupling derived
from the simultaneous BEMT--momentum balance (\S\ref{sec:soft}) is
\begin{empheq}[box=\fbox]{equation}
K_2 k_v \big|_{\text{self-consistent}} \;=\; \frac{\sigma a}{16\,\lambda_h},
\qquad \lambda_h = \sqrt{C_T/2}.
\label{eq:K2kv_consistent}
\end{empheq}
This is the formula to use in design.

% =====================================================================
\section{Inflow Softening and the Nonlinear Lead}
% =====================================================================
\label{sec:soft}

\subsection{What Inflow Softening Means}

The DC gain \(K_1/(1+K_2k_v)\) varies with the trim point because
\(k_v\propto 1/\bvi \propto 1/\sqrt{\bT}\). The factor \((1+K_2k_v)\) is the
\emph{inflow softening factor}---softening is strongest at low thrust.

\subsection{BEMT Derivation of the Softening Factor}

In hover, the coupled equations are
\begin{align}
C_T &= \frac{\sigma a}{2}\!\left(\frac{\thetat}{3}-\frac{\lambda}{2}\right), \tag{BEMT}\\
\lambda &= \sqrt{C_T/2}. \tag{Momentum}
\end{align}
Eliminating \(C_T\):
\begin{equation}
\lambda^2 = \frac{\sigma a}{4}\!\left(\frac{\thetat}{3}-\frac{\lambda}{2}\right),
\end{equation}
with closed form
\begin{equation}
\lambda(\thetat) = \frac{\sigma a}{16}\!\left[\sqrt{1+\frac{64\thetat}{3\sigma a}}-1\right].
\label{eq:lambda_of_theta}
\end{equation}

\paragraph{Trim sensitivity.} Differentiating implicitly:
\begin{equation}
2\lambda\,\frac{d\lambda}{d\thetat} = \frac{\sigma a}{4}\!\left(\frac{1}{3}-\frac{1}{2}\frac{d\lambda}{d\thetat}\right)
\;\;\Longrightarrow\;\;
\frac{d\lambda}{d\thetat} = \frac{\sigma a}{24\lambda + 3\sigma a/2}.
\end{equation}
Then
\begin{empheq}[box=\fbox]{equation}
\frac{dC_T}{d\thetat}\bigg|_{\text{steady}}
= \frac{\sigma a}{6}\cdot\underbrace{\frac{1}{1+\dfrac{\sigma a}{16\lambda}}}_{S(\lambda)}.
\label{eq:dCT_dtheta}
\end{empheq}
The factor \(\sigma a/6\) is the instantaneous (high-frequency) gain;
\(S(\lambda)\) is the softening factor, confirming~\eqref{eq:K2kv_consistent}.

\subsection{Magnitude Across the Envelope}

For a 600-class tail rotor (\(\sigma=0.138\), \(a=5.7\), hover \(\bCT=0.024\),
\(\blambda=0.110\)):

\begin{center}
\renewcommand{\arraystretch}{1.15}
\begin{tabular}{lcccc}
\toprule
Thrust & \(C_T\) & \(\lambda\) & \(K_2k_v=\sigma a/(16\lambda)\) & \(S(\lambda)\) \\
\midrule
25\% hover  & 0.006 & 0.055 & 0.89 & 0.53 \\
50\% hover  & 0.012 & 0.078 & 0.63 & 0.61 \\
Hover       & 0.024 & 0.110 & 0.45 & 0.69 \\
150\% hover & 0.036 & 0.135 & 0.36 & 0.73 \\
200\% hover & 0.048 & 0.156 & 0.32 & 0.76 \\
\bottomrule
\end{tabular}
\end{center}

The DC gain varies by roughly \(40\%\) across the operating envelope. A
fixed-gain lead compensator will mis-track by this amount at the extremes.

\subsection{Time Constant Varies Too}

The Pitt--Peters time constant
\(\tau_i = 0.849/(4\OmegaT\lambda)\) scales as \(1/\lambda\propto 1/\sqrt{C_T}\).
For the same 600-class rotor:

\begin{center}
\renewcommand{\arraystretch}{1.15}
\begin{tabular}{lc}
\toprule
Thrust & \(\tau_i\) [ms] \\
\midrule
25\% hover  & 5.5 \\
Hover       & 2.7 \\
200\% hover & 1.9 \\
\bottomrule
\end{tabular}
\end{center}

Both gain and time constant of the rotor depend on the trim operating point.
Both must be scheduled.

% =====================================================================
\section{Nonlinear Lead Compensator}
% =====================================================================

\subsection{Approach A: Gain Scheduling on Commanded Thrust (LPV)}

Schedule the lead-compensator coefficients on the current trim \(\lcmd\)
computed from the commanded thrust:
\begin{equation}
\lcmd(t) = \sqrt{\frac{\CTcmd(t)}{2}}
        = \sqrt{\frac{\Tcmd(t)}{2\rho\At(\OmegaT\Rt)^2}}.
\end{equation}
The scheduled compensator is
\begin{empheq}[box=\fbox]{equation}
\begin{aligned}
\tau_i(t) &= \frac{0.849}{4\,\OmegaT\,\lcmd(t)},\\[4pt]
\beta(t) &\equiv 1+K_2k_v = 1 + \frac{\sigma a}{16\,\lcmd(t)},\\[4pt]
\frac{\thetat(s)}{\Tt(s)} &= \frac{\tau_i(t)\,s+\beta(t)}{K_1\bigl(\tau_i(t)\,s+1\bigr)}.
\end{aligned}
\end{empheq}
In continuous state form, with a single internal state \(x\),
\begin{align}
\dot{x} &= -\frac{1}{\tau_i(t)}\,x + \frac{\beta(t)-1}{\tau_i(t)\,K_1}\,\Tcmd,\\
\thetat &= \frac{1}{K_1}\,\Tcmd + x.
\end{align}
The \(1/K_1\) feedforward gives the instantaneous inversion; the state \(x\)
adds the low-frequency correction.

\subsection{Approach B: Exact Static Inversion + Inflow Filter}

Bypass linearisation. The steady-state map \(\thetat\leftrightarrow C_T\)
is exact in closed form: from~\eqref{eq:CT_BEMT} with \(\lambda=\sqrt{C_T/2}\),
\begin{empheq}[box=\fbox]{equation}
\thetat^{\text{ss,cmd}} \;=\; \frac{6\,\CTcmd}{\sigma a} + \frac{3\,\lcmd}{2},
\qquad \lcmd = \sqrt{\CTcmd/2}.
\label{eq:static_inverse}
\end{empheq}
This static inverse, applied directly as feedforward, intentionally over-drives
the rotor for the duration of the inflow lag: the rotor produces a brief thrust
spike that builds inflow up to \(\lcmd\), after which thrust settles to \(\CTcmd\).
This is the nonlinear analogue of the lead compensator, and is structurally
simpler than the LPV form: a single algebraic expression in \(\CTcmd\).

For completeness, an inflow state with scheduled time constant can be
propagated alongside,
\begin{equation}
\dot{\lambda} = \frac{1}{\tau_i(\lcmd)}\bigl(\lcmd-\lambda\bigr),
\end{equation}
to expose the actual inflow as a state---useful for simulation, monitoring, or
hybridising with feedback.

\subsection{Approach C: Two-Point Approximation}

For embedded implementations where LPV is too heavy, interpolate the
compensator zero \(\beta/\tau_i\) and pole \(1/\tau_i\) linearly in
\(\sqrt{\Tcmd}\) between two design points (low and high thrust).

\subsection{Comparison}

\begin{center}
\renewcommand{\arraystretch}{1.2}
\begin{tabular}{p{4cm}p{3cm}p{3cm}p{4cm}}
\toprule
Compensator & Initial \(\thetat\) & Settled \(\thetat\) & Notes \\
\midrule
None (\(1/K_1\) only) & correct & undershoots ~31\% & thrust slow to build \\
Fixed lead at hover & correct & DC tracks at hover & overshoots at high thrust \\
LPV lead (A) & correct & DC tracks everywhere & needs schedule \\
Exact static inverse (B) & over-drives by \(\beta(\lcmd)\) & DC tracks exactly & cleanest, single expr.\ \\
\bottomrule
\end{tabular}
\end{center}

The static-inverse approach (B) is generally preferred: the nonlinearity falls
out of the BEMT algebra without a schedule, and the dynamic correction is a
single first-order state. The LPV form (A) is more rigorous for control analysis
but rarely justifies its complexity in a tail-rotor inner loop.

\subsection{Scheduling Variable Choice}

Should the schedule key on commanded thrust \(\Tcmd\) or on the actual
(estimated) thrust \(\Tt\)? Both have drawbacks. Scheduling on the command
assumes the rotor is tracking; scheduling on the measured thrust is
self-consistent but susceptible to noise. A practical compromise is to schedule
on a low-pass filtered command, with cutoff near \(1/\bar{\tau}_i\)---this is
``scheduling on slow variables'' in LPV terminology and is sound provided the
scheduling-variable bandwidth lies below the loop bandwidth.

% =====================================================================
\section{Parameter Ranges for RC Tail Rotors}
% =====================================================================

\subsection{Geometric and Operating Parameters}

\begin{center}
\renewcommand{\arraystretch}{1.15}
\begin{tabular}{llcccc}
\toprule
Symbol & Quantity & Micro & 450-class & 550--700 & 800+ \\
\midrule
\(\Rt\) & Tail radius [m]      & 0.03--0.05  & 0.06--0.08  & 0.10--0.13  & 0.14--0.17 \\
\(c\)    & Blade chord [m]       & 0.008--0.012 & 0.015--0.020 & 0.022--0.028 & 0.028--0.035 \\
\(N_b\)  & Number of blades       & 2 & 2 & 2 (occ.\ 4) & 2 or 4 \\
\(\OmegaT\) & Rotor speed [\si{\radian\per\second}] & 800--1500 & 600--1000 & 500--800 & 400--700 \\
\(\OmegaT\Rt\) & Tip speed [m/s] & 40--60 & 50--75 & 70--100 & 80--110 \\
\(\sigma\) & Solidity              & 0.10--0.16 & 0.12--0.18 & 0.12--0.18 & 0.12--0.20 \\
\(a\)    & Lift-curve slope [1/rad] & 5.5--5.75 & 5.5--5.75 & 5.5--5.75 & 5.5--5.75 \\
\bottomrule
\end{tabular}
\end{center}

\noindent The lift-curve slope sits near \(2\pi\approx 6.28\) in theory,
reduced to roughly 5.7 by finite-aspect-ratio and tip-loss effects, and falls
further at very low Reynolds (\(Re<\num{50000}\)).

\subsection{Aerodynamic Trim Values}

\begin{center}
\renewcommand{\arraystretch}{1.15}
\begin{tabular}{llcccc}
\toprule
Symbol & Quantity & Micro & 450 & 550--700 & 800+ \\
\midrule
\(\bT\)      & Hover thrust [N]              & 0.3--1 & 1--4 & 4--12 & 12--30 \\
\(\bvi\)     & Hover induced vel.\ [m/s]     & 4--8 & 5--10 & 7--13 & 9--15 \\
\(\btheta\)  & Trim pitch [deg]              & 6--10 & 6--10 & 6--10 & 6--10 \\
\(\tau_i\)   & Inflow time constant [ms]     & 4--10 & 8--20 & 15--35 & 25--60 \\
\bottomrule
\end{tabular}
\end{center}

\subsection{Derived Coefficients}

Using \(\rho = \SI{1.225}{\kilogram\per\cubic\metre}\):

\begin{center}
\renewcommand{\arraystretch}{1.15}
\begin{tabular}{llcccc}
\toprule
Symbol & Quantity & Micro & 450 & 550--700 & 800+ \\
\midrule
\(\At\)      & Disc area [\si{\square\metre}]   & 0.003--0.008 & 0.011--0.020 & 0.031--0.053 & 0.062--0.091 \\
\(K_1\)      & Pitch gain [N/rad]              & 3--10  & 10--40  & 40--150 & 150--400 \\
\(K_2\)      & Inflow gain [N\,s/m]            & 0.3--1.0 & 0.8--2.5 & 2.5--8 & 8--20 \\
\(k_v\)      & Inflow sensitivity [m/(s\,N)]   & 0.5--2.0 & 0.3--1.0 & 0.2--0.6 & 0.15--0.4 \\
\(K_2 k_v\)  & DC softening factor            & 0.4--1.0 & 0.4--1.0 & 0.4--1.0 & 0.4--1.0 \\
\bottomrule
\end{tabular}
\end{center}

\noindent The dimensionless product \(K_2 k_v\) is roughly the same across
classes---a physical-similarity result, since both factors scale with similar
combinations of \(\rho\At\), \(\OmegaT\Rt\), and \(\bvi\).

\subsection{Worked Example: 600-class}

\begin{align*}
\Rt &= 0.115~\text{m}, & c &= 0.025~\text{m}, & N_b &= 2,\\
\OmegaT &= 700~\text{rad/s}, & \OmegaT\Rt &= 80.5~\text{m/s},\\
\sigma &= \tfrac{2\cdot 0.025}{\pi\cdot 0.115} = 0.138, & a &= 5.7,\\
\bT &= 8~\text{N}, & \At &= 0.0415~\text{m}^2.
\end{align*}
Hence
\begin{align*}
\bvi &= \sqrt{8/(2\cdot 1.225\cdot 0.0415)} = 8.9~\text{m/s},\\
\tau_i &= 0.849\cdot 0.115/(4\cdot 8.9) \approx 2.7~\text{ms (Pitt--Peters)},\\
K_1 &= \tfrac{1.225\cdot 0.0415\cdot 80.5^2\cdot 0.138\cdot 5.7}{6} = 43~\text{N/rad},\\
K_2 &= \tfrac{1.225\cdot 0.0415\cdot 80.5\cdot 0.138\cdot 5.7}{4} = 8.1~\text{N\,s/m},\\
k_v &= 1/(4\cdot 1.225\cdot 0.0415\cdot 8.9) = 0.55~\text{m/(s\,N)},\\
\lambda_h &= \sqrt{0.024/2} = 0.110,\\
K_2k_v\big|_\text{consistent} &= \frac{\sigma a}{16\lambda_h} = \frac{0.787}{1.76} = 0.45.
\end{align*}

\subsection{Caveats}

Pitt--Peters time constants are systematically low compared to flight-test
data on full-scale rotors (factor of 2--4 underestimate); RC practitioners
often use \(\tau_i\approx 0.1\)--\(0.3\) rotor revolutions empirically.
Blade lift-curve slope drops noticeably at low Reynolds numbers
(\(Re\sim\num{30000}\)--\num{80000}) typical of RC tail-rotor tips,
especially with thin symmetric sections. Solidity calculations assume
rectangular blades; tapered or paddle blades need an equivalent solidity
weighted by \(r^2 c(r)\). Finally, the tail rotor sits in the main-rotor
wake, so \(V_\perp\) in~\eqref{eq:momentum} is dominated by main-rotor
inflow---the main source of tail-rotor effectiveness loss in turns.

% =====================================================================
\section{Summary}
% =====================================================================

The tail-rotor inner loop reduces to:

\begin{enumerate}[leftmargin=*]
\item A near-linear blade-element thrust law,
      \(\Tt = K_1\thetat - K_2\vi\), with constants from~\eqref{eq:K1K2}.
\item A first-order inflow state,
      \(\tau_i\dot{\vi}+\vi=\viss(\Tt)\), with \(\tau_i\) from~\eqref{eq:tau_i}.
\item A coupled steady-state map \(\thetat\leftrightarrow C_T\) with closed
      form~\eqref{eq:lambda_of_theta}.
\item A self-consistent DC softening factor \(\sigma a/(16\lambda_h)\) that
      replaces the naive product \(K_2 k_v\).
\item A nonlinear feedforward inverse---either the LPV lead
      \eqref{eq:comp_tf} scheduled on \(\lcmd\), or the exact static
      inverse~\eqref{eq:static_inverse}---that compensates for inflow
      softening across the full thrust envelope.
\end{enumerate}

\end{document}
